\subsection{Joint Distributions, Multivariate Laws, and Transformations}

\begin{definition} \label{def:joint-distributions-and-marginals} \hlt{Joint and Marginal Distributions}
\begin{enumerate}[label=(\roman*)]
\setlength{\itemsep}{0pt}
\item The \hlt{joint mass function} of discrete $X,Y$ is $f_{X,Y}(x,y)\defeq\mathbb{P}(X=x,Y=y)$.
\item The \hlt{discrete marginals} are $f_X(x)=\sum_y f_{X,Y}(x,y)$\\
and $f_Y(y)=\sum_x f_{X,Y}(x,y)$.
\item If $(X,Y)$ has joint density $f_{X,Y}$, the \hlt{continuous marginals} are $f_X(x)=\int_{-\infty}^{\infty}f_{X,Y}(x,y)\,dy$\\
and $f_Y(y)=\int_{-\infty}^{\infty}f_{X,Y}(x,y)\,dx$.
\item The \hlt{joint CDF} is $F_{X,Y}(x,y)\defeq\mathbb{P}(X\leq x,Y\leq y)$.
\end{enumerate}
\end{definition}

\begin{definition} \label{def:independent-random-variables-and-iid-samples} \hlt{Independent Random Variables and IID Samples}
\begin{enumerate}[label=(\roman*)]
\setlength{\itemsep}{0pt}
\item Random variables $X_1,\ldots,X_n$ are \hlt{independent}\\
if $\mathbb{P}(X_1\in B_1,\ldots,X_n\in B_n)=\prod_{i=1}^{n}\mathbb{P}(X_i\in B_i)$ for all $B_i\in\mathcal{B}(\R)$.
\item They are \hlt{identically distributed} if $\mathcal{L}(X_1)=\cdots=\mathcal{L}(X_n)$.
\item They are \hlt{iid} if they are independent and identically distributed.
\item The \hlt{random vector} $X=(X_1,\ldots,X_n)$ is a \hlt{random sample} when $X_1,\ldots,X_n$ are iid from a common law.
\end{enumerate}
\end{definition}

\begin{proposition} \label{prop:factorization-tests-for-random-variable-independence} \hlt{Factorization Tests for Independence}
\begin{enumerate}[label=(\roman*)]
\setlength{\itemsep}{0pt}
\item \hlt{Discrete factorization}: discrete $X,Y$ are independent iff $f_{X,Y}(x,y)=f_X(x)f_Y(y)$ for all $x,y$.
\item \hlt{Density factorization}: if $(X,Y)$ has joint density $f_{X,Y}$ and marginals $f_X,f_Y$,\\
then $X,Y$ are independent iff $f_{X,Y}(x,y)=f_X(x)f_Y(y)$ a.e.
\item \hlt{Rectangular support test}: if $f_{X,Y}(x,y)=g(x)h(y)\mathbf{1}_{A\times B}(x,y)$ for measurable $A,B$,\\
then $X,Y$ are independent after normalization.
\end{enumerate}
\begin{proof}
\begin{enumerate}[label=(\roman*)]
\setlength{\itemsep}{0pt}
\item[$(i)$] ($\Rightarrow$) If $X\perp\!\!\!\perp Y$, then $f_{X,Y}(x,y)=\mathbb{P}(X=x,Y=y)=\mathbb{P}(X=x)\mathbb{P}(Y=y)=f_X(x)f_Y(y)$.
\item[$(i)$] ($\Leftarrow$) If $f_{X,Y}=f_Xf_Y$, then for $A,B$, $\mathbb{P}(X\in A,Y\in B)=\sum_{x\in A}\sum_{y\in B}f_X(x)f_Y(y)=\mathbb{P}(X\in A)\mathbb{P}(Y\in B)$.
\item[$(ii)$] ($\Rightarrow$) If $X\perp\!\!\!\perp Y$, the product law has density $f_Xf_Y$, so uniqueness of densities gives $f_{X,Y}=f_Xf_Y$ a.e.
\item[$(ii)$] ($\Leftarrow$) If $f_{X,Y}=f_Xf_Y$ a.e., then for Borel $A,B$,\\
$\mathbb{P}(X\in A,Y\in B)=\int_A\int_Bf_X(x)f_Y(y)\,dy\,dx=\mathbb{P}(X\in A)\mathbb{P}(Y\in B)$.
\item[$(iii)$] Let $G=\int_Ag$ and $H=\int_Bh$. Normalization gives $GH=1$, while $f_X(x)=g(x)H\mathbf{1}_A(x)$ and\\
$f_Y(y)=h(y)G\mathbf{1}_B(y)$, hence $f_X(x)f_Y(y)=f_{X,Y}(x,y)$ a.e.; apply (ii).
\end{enumerate}
\hlt{Structural Commentary}: Independence is product structure of the joint law.\\
\hlt{Key Insight}: Independence means the joint law factors.
\end{proof}
\end{proposition}

\begin{example} \label{ex:bivariate-density-nonrectangular-region} \hlt{A Non-Rectangular Bivariate Density}\\
Let $S=\{(x,y):x^2\leq y\leq1\}$ and $f_{X,Y}(x,y)=cx^2y\mathbf{1}_S(x,y)$.\\
Since $S=\{(x,y):-1\leq x\leq1,\ x^2\leq y\leq1\}$, $1=c\int_{-1}^{1}\int_{x^2}^{1}x^2y\,dy\,dx=4c/21$, so $c=21/4$.\\
For $E=\{X\geq Y\}$, $S\cap E=\{(x,y):0\leq x\leq1,\ x^2\leq y\leq x\}$, hence $\mathbb{P}(E)=(21/4)\int_0^1\int_{x^2}^{x}x^2y\,dy\,dx=3/20$.\\
\hlt{Key Insight}: In bivariate integration, pick one variable first and let the other range over the corresponding slice.
\end{example}

\begin{figure}[H]
  \centering
  % Source integration: simplified redraw of the non-rectangular bivariate-density
% support and the event region from the attached random-variable chapter screenshots.
\begin{tikzpicture}[
  x=3.0cm,
  y=3.0cm,
  >=Stealth,
  line cap=round,
  line join=round,
  every node/.style={font=\scriptsize},
  axis/.style={->, gray!70, line width=0.45pt},
  tick/.style={gray!70, line width=0.4pt},
  support/.style={fill=gray!12, draw=none},
  event/.style={fill=blue!22, draw=none},
  boundary/.style={black!72, line width=0.75pt},
  event boundary/.style={blue!70!black, line width=0.9pt}
]
  \path[support, domain=-1:1, variable=\x]
    plot ({\x},{\x*\x}) -- (1,1) -- (-1,1) -- cycle;
  \path[event] (0,0) -- plot[domain=0:1, variable=\x] ({\x},{\x*\x}) -- cycle;

  \draw[axis] (-1.18,0) -- (1.22,0) node[black!78,right] {$x$};
  \draw[axis] (0,-0.08) -- (0,1.18) node[black!78,above] {$y$};

  \draw[boundary, domain=-1.04:1.04, samples=120, variable=\x] plot ({\x},{\x*\x});
  \draw[boundary] (-1,1) -- (1,1);
  \draw[event boundary] (0,0) -- (1,1);

  \foreach \x/\lab in {-1/-1,0/0,1/1}
    \draw[tick] (\x,0.018) -- (\x,-0.018) node[black!78,below] {$\lab$};
  \node[black!78,anchor=south east] at (-0.05,1.02) {$1$};
\end{tikzpicture}

  \caption{\textcolor{gray!70!black}{Gray}: support $x^2\leq y\leq1$; \textcolor{blue!70!black}{blue}: event $X\geq Y$ ($x^2\leq y\leq x$).}
  \label{fig:parabolic-bivariate-region}
\end{figure}

\begin{definition} \label{def:multinomial-and-multivariate-normal} \hlt{Two Multivariate Families}
\begin{enumerate}[label=(\roman*)]
\setlength{\itemsep}{0pt}
\item \hlt{Multinomial}: for $n\in\N$ and $p\in[0,1]^k$ with $\sum_{j=1}^{k}p_j=1$, $X\sim\operatorname{Multinomial}(n,p)$\\
if $\mathbb{P}(X=x)=n!(x_1!\cdots x_k!)^{-1}\prod_{j=1}^{k}p_j^{x_j}$ for $x\in\N_0^k$ with $\sum_{j=1}^{k}x_j=n$; then $X_j\sim\operatorname{Binomial}(n,p_j)$.
\item \hlt{Multivariate normal}: for $\mu\in\R^k$ and symmetric positive definite $\Sigma\in\R^{k\times k}$, $X\sim N(\mu,\Sigma)$\\
if $f_X(x)=((2\pi)^k\det\Sigma)^{-1/2}\exp(-(x-\mu)^t\Sigma^{-1}(x-\mu)/2)$ for $x\in\R^k$.\\
$N(0,I_k)$ has independent $N(0,1)$ coordinates.
\end{enumerate}
\end{definition}

\begin{proposition} \label{prop:multivariate-normal-affine-standardization} \hlt{Multivariate Normal Affine Standardization}\\
Let $\mu\in\R^k$ and $\Sigma\in\R^{k\times k}$ with $\Sigma\succ0$. For a random vector $X\in\R^k$, $X\sim N(\mu,\Sigma)$ iff $Z\defeq\Sigma^{-1/2}(X-\mu)\sim N(0,I_k)$; equivalently, if $Z\sim N(0,I_k)$, then $\mu+\Sigma^{1/2}Z\sim N(\mu,\Sigma)$.
\begin{proof}
\begin{enumerate}[label=\roman*.]
\setlength{\itemsep}{0pt}
\item[$(\Rightarrow)$] Set $Z=\Sigma^{-1/2}(X-\mu)$, so $x=\mu+\Sigma^{1/2}z$ and $|\det\Sigma^{1/2}|=(\det\Sigma)^{1/2}$.\\
By change of variables, $f_Z(z)=f_X(\mu+\Sigma^{1/2}z)(\det\Sigma)^{1/2}=(2\pi)^{-k/2}\exp(-z^tz/2)$, hence $Z\sim N(0,I_k)$.
\item[$(\Leftarrow)$] Let $X=\mu+\Sigma^{1/2}Z$ with $Z\sim N(0,I_k)$. Since $z=\Sigma^{-1/2}(x-\mu)$ and $|\det\Sigma^{-1/2}|=(\det\Sigma)^{-1/2}$,\\
$f_X(x)=f_Z(\Sigma^{-1/2}(x-\mu))(\det\Sigma)^{-1/2}=((2\pi)^k\det\Sigma)^{-1/2}\exp(-(x-\mu)^t\Sigma^{-1}(x-\mu)/2)$, so $X\sim N(\mu,\Sigma)$.
\end{enumerate}
\hlt{Structural Commentary}: Positive definite covariance is an affine deformation of the standard normal.\\
\hlt{Key Insight}: Standardize by $\Sigma^{-1/2}$; reconstruct by $\Sigma^{1/2}$.
\end{proof}
\end{proposition}

\begin{proposition} \label{prop:multivariate-normal-basic-calculus} \hlt{Basic Multivariate Normal Calculus}\\
Let $X=(X_a,X_b)\sim N(\mu,\Sigma)$ with $X_a\in\R^r$, $X_b\in\R^s$, $\mu=(\mu_a,\mu_b)$, and $\Sigma=\begin{pmatrix}\Sigma_{aa}&\Sigma_{ab}\\ \Sigma_{ba}&\Sigma_{bb}\end{pmatrix}\succ0$.
\begin{enumerate}[label=(\roman*)]
\setlength{\itemsep}{0pt}
\item \hlt{Marginal}: $X_a\sim N(\mu_a,\Sigma_{aa})$.
\item \hlt{Conditional}: $X_b\mid X_a=x_a\sim N(\mu_b+\Sigma_{ba}\Sigma_{aa}^{-1}(x_a-\mu_a),\Sigma_{bb}-\Sigma_{ba}\Sigma_{aa}^{-1}\Sigma_{ab})$.
\item \hlt{Linear form}: for any vector $a$, $a^tX\sim N(a^t\mu,a^t\Sigma a)$.
\end{enumerate}
\begin{proof}
\begin{enumerate}[label=(\roman*)]
\setlength{\itemsep}{0pt}
\item Let $P_aX=X_a$. By Proposition \ref{prop:multivariate-normal-affine-standardization}, $X=\mu+\Sigma^{1/2}Z$ with $Z\sim N(0,I)$.\\
Hence $X_a=P_a\mu+P_a\Sigma^{1/2}Z\sim N(\mu_a,P_a\Sigma P_a^t)=N(\mu_a,\Sigma_{aa})$.
\item Set $S\defeq\Sigma_{bb}-\Sigma_{ba}\Sigma_{aa}^{-1}\Sigma_{ab}$ and $m(x_a)\defeq\mu_b+\Sigma_{ba}\Sigma_{aa}^{-1}(x_a-\mu_a)$.\\
The block determinant and inverse identities give $\det\Sigma=\det\Sigma_{aa}\det S$ and\\
$(x-\mu)^t\Sigma^{-1}(x-\mu)=(x_a-\mu_a)^t\Sigma_{aa}^{-1}(x_a-\mu_a)+(x_b-m(x_a))^tS^{-1}(x_b-m(x_a))$.\\
Thus $f_{b\mid a}(x_b\mid x_a)=f_{a,b}(x_a,x_b)/f_a(x_a)$ is the $N(m(x_a),S)$ density.
\item With $X=\mu+\Sigma^{1/2}Z$, $a^tX=a^t\mu+(\Sigma^{1/2}a)^tZ$.\\
Since $Z$ has independent $N(0,1)$ coordinates, $(\Sigma^{1/2}a)^tZ\sim N(0,a^t\Sigma a)$, hence $a^tX\sim N(a^t\mu,a^t\Sigma a)$.
\end{enumerate}
\hlt{Structural Commentary}: Marginals, conditionals, and linear projections of Gaussian vectors remain Gaussian.\\
\hlt{Key Insight}: Gaussian calculus is affine calculus plus Schur complements.
\end{proof}
\end{proposition}

\begin{method} \label{meth:one-variable-transformation-cdf-method} \hlt{One-Variable Transformations}\\
Let $Y=r(X)$.
\begin{enumerate}[label=(\roman*)]
\setlength{\itemsep}{0pt}
\item \hlt{Discrete pushforward}: if $X$ is discrete, then $f_Y(y)=\mathbb{P}(Y=y)=\mathbb{P}(X\in r^{-1}(\{y\}))=\sum_{x:r(x)=y}f_X(x)$.
\item \hlt{CDF method}: if $X$ has density $f_X$, set $A_y\defeq\{x:r(x)\leq y\}$, compute $F_Y(y)=\mathbb{P}(Y\leq y)=\int_{A_y}f_X(x)\,dx$, and use $f_Y(y)=F_Y'(y)$ where differentiable.
\item \hlt{Monotone change of variables}: if $r$ is strictly monotone with inverse $s=r^{-1}$,\\
then $f_Y(y)=f_X(s(y))|s'(y)|$ on $r(\operatorname{supp}X)$.
\end{enumerate}
\hlt{Key Insight}: The CDF method works because it translates the transformed event back into an event about $X$.
\end{method}

\begin{example} \label{ex:log-transform-exponential} \hlt{Log Transform of an Exponential Variable}\\
Let $X\sim\operatorname{Exp}(1)$ and $Y=\log X$. Since $r(x)=\log x$ has inverse $s(y)=e^y$ with $s'(y)=e^y$, the change-of-variables formula gives $f_Y(y)=f_X(e^y)e^y=e^ye^{-e^y}$ for $y\in\R$.
\end{example}

\begin{example} \label{ex:square-transform-nonmonotone-uniform} \hlt{A Nonmonotone Square Transform}\\
Let $X\sim\operatorname{Uniform}(-1,3)$ and $Y=X^2$. Since $f_X=1/4$ on $[-1,3]$, $r^{-1}(\{y\})\cap[-1,3]=\{-\sqrt y,\sqrt y\}$ for $0<y<1$ and $\{\sqrt y\}$ for $1<y<9$. Hence $f_Y(y)=1/(4\sqrt y)\mathbf{1}_{(0,1)}(y)+1/(8\sqrt y)\mathbf{1}_{(1,9)}(y)$.\\
\hlt{Key Insight}: Non-one-to-one transformations require splitting the range into cases.
\end{example}

\begin{method} \label{meth:several-variable-transformation-cdf-method} \hlt{Several-Variable Transformations}\\
For $Z=r(X,Y)$ with joint density $f_{X,Y}$:
\begin{enumerate}[label=(\roman*)]
\setlength{\itemsep}{0pt}
\item \hlt{Sublevel set}: define $A_z\defeq\{(x,y):r(x,y)\leq z\}$.
\item \hlt{CDF integral}: compute $F_Z(z)=\mathbb{P}(Z\leq z)=\iint_{A_z}f_{X,Y}(x,y)\,dx\,dy$.
\item \hlt{Density}: differentiate where possible, $f_Z(z)=F_Z'(z)$.
\end{enumerate}
\end{method}

\begin{example} \label{ex:sum-of-two-independent-uniforms} \hlt{Sum of Two Independent Uniforms}\\
Let $X_1,X_2\sim\operatorname{Uniform}(0,1)$ be independent and $Y=X_1+X_2$.\\
Since $f_{X_1,X_2}=\mathbf{1}_{[0,1]^2}$, $F_Y(y)$ is the area of $A_y=\{(x_1,x_2)\in[0,1]^2:x_1+x_2\leq y\}$,\\
so $F_Y(y)=(y^2/2)\mathbf{1}_{[0,1]}(y)+(1-(2-y)^2/2)\mathbf{1}_{(1,2]}(y)+\mathbf{1}_{(2,\infty)}(y)$. Hence $f_Y(y)=y\mathbf{1}_{(0,1)}(y)+(2-y)\mathbf{1}_{(1,2)}(y)$.
\end{example}
